\section{Knowledge-Map Reconstruction Agent}
\label{sec:reconstruction-agent}

The benchmark makes active user modeling measurable, but it does not by itself determine how an agent
should reconstruct the hidden map.

The central method question is how to turn sparse, open-ended answers into an explicit state estimate
and how to choose the next question from that estimate.

We propose \emph{MapProbe}, an evidence--belief--probe control loop that separates these decisions into
inspectable objects.

MapProbe is the full research design. The experimental Evidence-Calibrated Diagnostic Agent (ECDA)
implements a narrower slice at the time of writing.

The implemented system already enforces the visibility boundary, initializes a full-graph working-map
shell, validates node updates, records visible support, and provides a simple prompt-based LLM agent.

ECDA additionally persists independent per-node mastery marginals, applies deterministic Bayesian-style
updates to model-proposed likelihoods, and selects deterministically among multiple model-proposed
question candidates.

The full direct/indirect evidence ledger, typed graph messages, target-before-language verifier, and
calibrated stopping rule described below are not yet implemented. No MapProbe or ECDA mechanism has
been comparatively or empirically validated.

We specify them now so that missing results leave a falsifiable method contract rather than being
replaced by benchmark-description prose.

\subsection{External Protocol and Current Baseline}
\label{sec:agent}

The agent protocol constrains what must be observable, not how a policy must reason internally.

Every tested agent maintains an assessment for each node,

\begin{equation}
    \hat{s}_{i,t} = (\hat{m}_{i,t}, c_{i,t}, n_{i,t}, R_{i,t}),
\end{equation}
where $\hat{m}_{i,t} \in \{\unknown,0,\ldots,5\}$ is assessed mastery, $c_{i,t}$ is confidence,
$n_{i,t}$ is an optional note, and $R_{i,t}$ contains supporting visible turns.

All nodes begin as \unknown. The schema is not claimed to be cognitively privileged; it makes the
update trajectory inspectable and keeps unsupported estimates scoreable.

The runtime exposes only benchmark-relevant operations: read the public graph or current working map,
update node assessments with visible support, ask one diagnostic question, and finalize the full map.

These operations expose two links: answers should change the working map, and that map should influence
later probes.

The runtime enforces turn order and visibility but performs no diagnosis for the tested agent.

The implemented simple LLM baseline uses this interface through two structured prompts.

One interprets the latest answer and emits hard node updates. The other emits a primary target, optional
connected targets, a mastery boundary, a selection reason, and one question.

It provides a strong schema-aware prompting control, but it has no explicit probabilistic uncertainty,
observation likelihood, formal graph propagation, or target-utility calculation.

ECDA adds an experimental observation-likelihood and node-marginal update path. It also validates at
least three complete question candidates and selects one with a deterministic utility over
model-estimated information gain, coverage, graph leverage, redundancy, complexity, and outcome-model
confidence. This is candidate-level selection after language generation, not the full MapProbe
target-first planner and verifier.

MapProbe is evaluated as a mechanism change relative to this baseline, not merely as a larger prompt.

\subsection{Design Objective and State}

At turn $t$, MapProbe maintains a belief object for every public node $v$:

\begin{equation}
    B_t(v) = \langle p_t(\ell\mid v), D_t(v), I_t(v), C_t(v) \rangle,
\end{equation}
where $p_t(\ell\mid v)$ is a distribution over mastery levels $\ell\in\{0,\ldots,5\}$,
$D_t(v)$ and $I_t(v)$ are direct and graph-inferred evidence ledgers, and $C_t(v)$ records
contradictions, ambiguity, and missing coverage.

This extends the benchmark's externally visible hard assessment without changing the required
submission schema. A submitted level is a projection of the belief state, not the entire state itself.

Sequential state maintenance is motivated by knowledge tracing, but the prediction target is different.

Knowledge-tracing models update a latent state to predict future item responses
\citep{piech2015dkt,liu2022pykt}; \knowact directly scores the reconstructed state over every graph
node.

Dynamic profiling similarly supports explicit profile updates, while showing why reconstruction
quality should be separated from a downstream response \citep{zhao2025rlpa}.

MapProbe therefore retains the belief and its evidence even when a later answer would be plausible
without it.

All priors are fixed before test interaction.

A uniform prior is the default zero-shot condition. An empirical prior estimated from permitted
training episodes is a separate baseline.

Test profiles, simulator traces, and scorer outputs never affect initialization. A node may remain
\unknown when evidence is insufficient, but its abstention threshold is frozen on development data.

\subsection{Rubric-Grounded Evidence Update}

After observing $(q_t,a_t)$, an evidence interpreter receives the visible turn, its diagnostic plan,
and the public definitions and rubrics of the targeted nodes.

It produces typed observations containing the implicated node, mastery boundary, polarity, strength,
visible turn identifier, and direct or indirect provenance.

It must distinguish demonstrated reasoning from self-reported confidence and preserve uncertainty for
incomplete, off-target, or internally inconsistent answers.

This decomposition follows the empirical separation between evidence acquisition, diagnosis, and
action in LongTutor \citep{li2026longtutor}.

Conversational diagnosis and assessment-driven tutoring also motivate persistent turn-level state
rather than a regenerated whole-dialogue judgment \citep{yao2026parld,li2026scaffoldlm}.

We do not infer from these works that multiple LLM agents are necessary. MapProbe uses a typed boundary
because it can be checked and ablated, not because role multiplicity is itself a contribution.

For a direct observation $o_t$ about node $v$, the belief update is

\begin{equation}
    p_t(\ell\mid v)
    \propto p_{t-1}(\ell\mid v) P(o_t\mid \ell,v).
\end{equation}
The current ECDA slice asks the model for a six-level observation-likelihood vector and applies the
update deterministically. The likelihoods are zero-shot estimates, not calibrated or learned
probabilities. A frozen ordinal table remains a reproducible baseline, and a learned observation model
becomes eligible only after adjudicated training interactions exist.

New evidence does not overwrite conflicting evidence. Contradiction increases uncertainty and remains a
candidate for follow-up.

Repeated semantically equivalent observations receive a redundancy discount, so verbosity is not
treated as independent support.

A bounded verifier checks node identity, rubric compatibility, visible provenance, and whether the cited
answer supports the claimed boundary.

Its value must be tested against the same interpreter without verification and against human evidence
annotations. Agreement between two prompts from the same model is not independent semantic validation.

\subsection{Graph-Aware Soft Inference}

The graph can make a probe informative beyond one node, but it can also amplify an authored error.

MapProbe therefore treats edges as evidence-routing priors rather than deterministic mastery rules. For
an edge $(u,r,v)$, it forms an attenuated relation-specific message:

\begin{equation}
    m_{u\rightarrow v}^{(r)}(\ell)
    = \alpha_r q_r(\ell\mid B_t(u)),
    \qquad 0\leq\alpha_r<1,
\end{equation}
Here, $q_r$ follows the declared semantics of relation $r$.

A mastered prerequisite may increase the plausibility of downstream mastery, but it cannot copy a
level. Negative evidence may motivate probing a prerequisite without proving downstream non-mastery.

Direct and propagated evidence remain separate, and the initial method allows only one propagation hop.

This restriction is substantive: multi-hop messages can create confident correlated errors when graph
edges or node granularity are imperfect.

Unlike sequential state update and adaptive item selection, typed propagation over an authored
diagnostic graph lacks a direct complete precedent in the focused literature.

It is a \knowact hypothesis that requires no-graph, untyped-graph, and deliberately perturbed-graph
controls before supporting a graph-aware diagnosis claim.

\subsection{Diagnostic Target Planning}

MapProbe selects a target before generating question language. For a node or small connected target set
$S$, the planner scores

\begin{align}
    U_t(S) ={}& \lambda_1 \widehat{\mathrm{IG}}_t(S)
    + \lambda_2 \operatorname{Coverage}_t(S)
    + \lambda_3 \operatorname{BoundaryRisk}_t(S) \nonumber\\
    &+ \lambda_4 \operatorname{GraphReach}_t(S)
    - \lambda_5 \operatorname{Redundancy}_t(S)
    - \lambda_6 \operatorname{Cost}_t(S).
\end{align}
Expected information gain estimates the reduction in full-map uncertainty under plausible answers.

Coverage rewards direct evidence for untouched nodes. Boundary risk prioritizes posterior mass around
a high-loss adjacent-level decision.

Graph reach is attenuated by relation type. Redundancy discounts repeated probes, and cost reflects
remaining turns and expected answer burden.

Adaptive testing establishes that question selection should be compared at matched lengths, not judged
only by the final test score \citep{ghosh2021bobcat,liu2025facd}.

Selective user-information acquisition further motivates reporting outcome, query cost, and an explicit
insufficiency decision \citep{wu2025raise}.

These results do not validate MapProbe's information-gain estimate: fixed items and user attributes are
simpler than open diagnostic questions. Maximum-entropy and entropy-plus-coverage remain controls.

The plan is written into the existing visible action schema: a primary node, optional secondary nodes,
the mastery boundary to discriminate, and a concise selection reason.

Exposing this object permits an independent check of whether later language and updates follow the
claimed decision.

\subsection{Question Realization, Stopping, and Submission}

The realizer converts the selected plan into one natural diagnostic question.

The question should elicit reasoning, application, transfer, comparison, or self-correction. It avoids
revealing rubrics or expected answers and does not bundle unrelated targets.

A verifier can accept, rewrite, or reject the question based on plan alignment and answerability. Its
extra model calls and correlated-judge risk are reported, not hidden inside the method.

The agent stops when the turn budget is exhausted or the best estimated marginal utility falls below a
frozen threshold.

Early stopping also requires high-priority coverage or explicit abstention. Remaining contradictions
must not be profitably separable within the available budget.

This decision is compared with fixed-turn policies at every budget. Fewer questions count as improved
efficiency only if reconstruction quality is preserved.

At finalization, MapProbe projects every node belief to $0$--$5$ or \unknown, a diagnostic-confidence
category, a note, and visible supporting turns.

The endpoint score remains deterministic. Selective-risk and calibration measures distinguish a
cautious agent from one that improves mean error with unsupported confident predictions.

\subsection{Failure Handling and Audit Trail}

MapProbe treats invalid intermediate output as an experimental event rather than silently repairing it
with hidden benchmark logic.

A structurally invalid evidence record is retried within a fixed call budget. If the retry fails, the turn
is an interpretation failure and the previous belief is preserved.

A weak observation may reduce no entropy. A rejected question receives one bounded rewrite and then
falls back to the next ranked target.

Provider failure, schema failure, verifier rejection, and voluntary abstention remain distinct trace
outcomes.

Every accepted update records the pre-update belief, cited visible turn, direct observation, likelihood
rule, post-update belief, and any outgoing graph message.

Every action records candidate scores, the selected target, rejected alternatives, generated question,
verifier decision, and remaining budget.

The post-episode trace never reveals hidden map or simulator state to the tested policy. It distinguishes
a poor observation model from a poor planner, plan, or linguistic realization.

The method also places explicit limits on self-correction. A verifier may reject an unsupported record,
but it may not invent evidence, query the hidden profile, or replace a node belief.

Graph propagation can change an indirect belief component but cannot add a direct supporting turn.
These restrictions keep the provenance claim mechanically testable.

The following claim contract determines how the agent may be described after experiments:

\begin{itemize}[leftmargin=1.5em]
    \item explicit state is beneficial only if the hard-label ablation is worse under matched calls;
    \item graph-aware diagnosis is beneficial only if no-graph and perturbed-graph controls support it;
    \item planning is informative only if it beats random, fixed, and maximum-entropy target policies;
    \item verification improves quality only if its ablation and human annotations agree; and
    \item stopping is efficient only if it reduces interaction without degrading reconstruction.
\end{itemize}

Absent the corresponding evidence, the paper will describe each item as an implemented mechanism, not
as an explanation for performance.

\subsection{Implementation and Claim Boundary}

\begin{table}[t]
\caption{Method status. ECDA implements a narrow experimental slice; unimplemented MapProbe
components retain paper space, and no component supports a benefit claim before matched comparison.}
\label{tab:agent-status}
\centering
\small
\begin{tabularx}{\linewidth}{>{\raggedright\arraybackslash}Xcc}
\toprule
\textbf{Component} & \textbf{Implemented} & \textbf{Empirically tested} \\
\midrule
Visibility boundary, full-map shell, validated hard updates & Yes & Unit/runtime checks only \\
Simple prompt-based LLM update and target-plan baseline & Yes & No comparative study \\
Checkpointed node-marginal belief and likelihood update & ECDA slice & No comparative study \\
Deterministic multi-candidate question utility & ECDA slice & No comparative study \\
Direct/indirect evidence ledger and verifier & No & No \\
Typed, attenuated graph messages & No & No \\
Target-before-language planner and question verifier & No & No \\
Calibrated stopping and selective submission analysis & No & No \\
\bottomrule
\end{tabularx}
\end{table}

Accordingly, this section contributes a concrete, testable agent design rather than evidence of improved
performance.

The next section fixes the comparisons that would justify stronger claims and the results that would
falsify them.

Future implementation detail may refine these modules, but it may not erase the explicit state,
acquisition, graph-use, and stopping hypotheses that define the method contribution.
